
\subsection{Especificación textual de los casos de uso}
La línea base identifica 40 casos de uso y los 40 se especifican mediante actor, precondición y escenario observable. Esta relación permite calcular M1b sin confundir un caso dibujado con un caso efectivamente especificado.

\begin{landscape}
\scriptsize
\begin{longtable}{|p{1.2cm}|p{1.9cm}|p{2.4cm}|p{5.0cm}|p{9.2cm}|}
\caption{Especificación textual resumida de CU-01 a CU-40.}\\
\hline\textbf{CU} & \textbf{RF} & \textbf{Actor} & \textbf{Precondición} & \textbf{Escenario principal / resultado observable}\\\hline
\endfirsthead
\hline\textbf{CU} & \textbf{RF} & \textbf{Actor} & \textbf{Precondición} & \textbf{Escenario principal / resultado observable}\\\hline
\endhead
CU-01 & RF-AUT-01 & Personal autorizado & Usuario registrado con credenciales activas. & Dado un usuario con credenciales válidas, cuando inicie sesión, entonces el sistema deberá autenticarlo y habilitar únicamente las opciones permitidas para su rol; cuando cierre sesión, deberá invalidar la sesión activa. \\ \hline
CU-02 & RF-AUT-02 & Administrador & Usuario autenticado y autorizado para el caso de uso. & Dado un usuario autenticado con un rol definido, cuando intente ejecutar una operación, entonces el sistema deberá permitirla solo si el permiso está asignado a ese rol y deberá denegarla en caso contrario. \\ \hline
CU-20 & RF-AUT-03 & Administrador & Usuario autenticado y autorizado para el caso de uso. & Dado un cambio crítico confirmado, cuando el sistema lo registre, entonces deberá conservar en auditoría el actor, fecha, acción, valor anterior y valor nuevo asociados al cambio. \\ \hline
CU-03 & RF-CLI-01 & Recepcionista & Usuario autenticado y autorizado para el caso de uso. & Dado un formulario con código interno, nombre de uso y al menos un medio de contacto válido, cuando Recepción confirme el registro, entonces el sistema deberá crear un único cliente y no deberá solicitar datos de salud, biometría, dirección domiciliaria ni imagen corporal. \\ \hline
CU-04 & RF-CLI-02 & Recepcionista & Usuario autenticado y autorizado para el caso de uso. & Dado un criterio de búsqueda válido, cuando Recepción consulte un cliente, entonces el sistema deberá devolver las coincidencias y mostrar únicamente membresía, último pago, asistencia y novedades autorizadas por el rol. \\ \hline
CU-05 & RF-CLI-03 & Recepcionista & Usuario autenticado y autorizado para el caso de uso. & Dado que existe un cliente coincidente por código o contacto, cuando Recepción intente registrar o reactivar el perfil, entonces el sistema deberá advertir la coincidencia y permitir actualizar o reactivar el registro existente sin crear un duplicado. \\ \hline
CU-06 & RF-MEM-01 & Administrador & Usuario autenticado y autorizado para el caso de uso. & Dado un usuario autorizado y los datos obligatorios de un plan o promoción, cuando guarde la configuración, entonces el sistema deberá conservar nombre, duración, precio, vigencia y estado y deberá mostrar el registro actualizado. \\ \hline
CU-07 & RF-MEM-02 & Recepcionista & Usuario autenticado y autorizado para el caso de uso. & Dado un plan vigente y un pago confirmado, cuando Recepción active o renueve la membresía, entonces el sistema deberá calcular y guardar las fechas de inicio y vencimiento según la duración del plan. \\ \hline
CU-08 & RF-MEM-03 & Personal autorizado & Usuario autenticado y autorizado para el caso de uso. & Dado un cliente con membresía registrada, cuando un usuario autorizado consulte su vigencia, entonces el sistema deberá mostrar estado, plan, fecha de inicio y fecha de vencimiento de la membresía seleccionada. \\ \hline
CU-09 & RF-MEM-04 & Recepcionista & Usuario autenticado y autorizado para el caso de uso. & Dado que existen membresías próximas a vencer o vencidas, cuando se ejecute la consulta diaria, entonces el sistema deberá separar las que vencen dentro de tres días de las ya vencidas y no deberá enviar mensajes externos de forma automática. \\ \hline
CU-10 & RF-PAG-01 & Recepcionista & Usuario autenticado y autorizado para el caso de uso. & Dado un pago con monto, concepto, medio, referencia y fecha válidos, cuando Recepción confirme el registro, entonces el sistema deberá guardar el pago y generar un comprobante interno vinculado al cliente. \\ \hline
CU-21 & RF-PAG-02 & Administrador & Usuario autenticado y autorizado para el caso de uso. & Dado un pago con estado pendiente, cuando un Administrador lo confirme, rechace o corrija, entonces el sistema deberá actualizar el estado y conservar el motivo y la evidencia de auditoría de la decisión. \\ \hline
CU-11 & RF-ASI-01 & Recepcionista & Usuario autenticado y autorizado para el caso de uso. & Dado un cliente identificado, cuando Recepción registre la asistencia por QR, entonces el sistema deberá guardar fecha, hora y método QR; si se utiliza contingencia manual, deberá exigir responsable y motivo. \\ \hline
CU-12 & RF-ASI-02 & Personal autorizado & Usuario autenticado y autorizado para el caso de uso. & Dado un filtro por cliente, fecha o turno, cuando un usuario autorizado consulte asistencias, entonces el sistema deberá mostrar únicamente los registros que cumplan el filtro y los conteos correspondientes. \\ \hline
CU-13 & RF-INV-01 & Administrador & Usuario autenticado y autorizado para el caso de uso. & Dado un producto con código, nombre, precio, unidad, stock mínimo y estado válidos, cuando el Administrador confirme el registro, entonces el sistema deberá guardar el producto y permitir su consulta posterior. \\ \hline
CU-14 & RF-INV-02 & Personal autorizado & Usuario autenticado y autorizado para el caso de uso. & Dado un producto existente y un movimiento válido, cuando un usuario autorizado registre una entrada o ajuste, entonces el sistema deberá actualizar el saldo y conservar cantidad, motivo, responsable y saldo resultante. \\ \hline
CU-15 & RF-VEN-01 & Recepcionista & Usuario autenticado y autorizado para el caso de uso. & Dado que todos los ítems de una venta tienen stock suficiente, cuando Recepción confirme la venta, entonces el sistema deberá registrar la venta y descontar las existencias en una única transacción; si falla una operación, no deberá aplicar cambios parciales. \\ \hline
CU-22 & RF-INV-03 & Administrador & Usuario autenticado y autorizado para el caso de uso. & Dado un período de consulta, cuando el Administrador solicite el control de stock, entonces el sistema deberá listar los productos con stock igual o inferior al mínimo y resumir entradas, salidas y ventas del período. \\ \hline
CU-16 & RF-CAJ-01 & Recepcionista & Usuario autenticado y autorizado para el caso de uso. & Dado un turno con pagos y ventas registrados, cuando Recepción ejecute el cierre, entonces el sistema deberá totalizar por medio de pago, comparar contra los valores declarados y registrar cualquier diferencia junto con la entrega del turno. \\ \hline
CU-17 & RF-NOV-01 & Personal autorizado & Usuario autenticado y autorizado para el caso de uso. & Dado un usuario autorizado y una novedad con categoría, detalle, responsable y vencimiento, cuando la registre, entonces el sistema deberá crearla con un estado válido y conservar los cambios posteriores de estado. \\ \hline
CU-23 & RF-INS-01 & Administrador & Usuario autenticado y autorizado para el caso de uso. & Dado un instructor registrado, cuando el Administrador actualice su operación, entonces el sistema deberá guardar turnos, capacidad, disponibilidad y novedades operativas sin solicitar datos de salud. \\ \hline
CU-18 & RF-RUT-01 & Instructor & Usuario autenticado y autorizado para el caso de uso. & Dado un cliente autorizado y un catálogo de ejercicios disponible, cuando el Instructor cree una rutina, entonces el sistema deberá guardar objetivo operativo, ejercicios, series, repeticiones, descanso y días, y deberá asignarla al cliente seleccionado. \\ \hline
CU-19 & RF-RUT-02 & Instructor & Usuario autenticado y autorizado para el caso de uso. & Dado una rutina existente, cuando el Instructor modifique ejercicios o parámetros y confirme el cambio, entonces el sistema deberá crear una nueva versión y mantener las versiones anteriores en solo lectura. \\ \hline
CU-24 & RF-RUT-03 & Instructor & Usuario autenticado y autorizado para el caso de uso. & Dado una rutina activa, cuando el Instructor registre seguimiento, entonces el sistema deberá guardar fecha, cumplimiento, repeticiones y observación operativa y deberá rechazar campos destinados a diagnósticos, lesiones, peso o medidas corporales reales. \\ \hline
CU-25 & RF-REP-01 & Administrador & Usuario autenticado y autorizado para el caso de uso. & Dado un período y filtros válidos, cuando el Administrador consulte el panel o reporte, entonces el sistema deberá calcular los indicadores con datos almacenados, mostrar los resultados filtrados y permitir exportarlos a CSV. \\ \hline
CU-26 & RF-MEM-05 & Cliente & Usuario autenticado y autorizado para el caso de uso. & Dado un Cliente autenticado, cuando consulte su membresía, entonces el sistema deberá mostrar únicamente el estado, plan y vigencia de su propia membresía y no deberá habilitar edición. \\ \hline
CU-27 & RF-PAG-03 & Cliente & Usuario autenticado y autorizado para el caso de uso. & Dado un Cliente autenticado, cuando consulte sus pagos, entonces el sistema deberá mostrar únicamente sus registros de pago y el estado de cada uno en modo de solo lectura. \\ \hline
CU-28 & RF-RUT-04 & Cliente & Usuario autenticado y autorizado para el caso de uso. & Dado un Cliente autenticado con una rutina asignada, cuando consulte su rutina, entonces el sistema deberá mostrar la versión activa, las versiones históricas y el seguimiento autorizado sin permitir modificar la rutina vigente. \\ \hline
CU-29 & RF-INV-04 & Dueño & Usuario autenticado y autorizado para el caso de uso. & Dado un usuario autorizado, cuando registre, edite o desactive una categoría, entonces el sistema deberá guardar el cambio y reflejar el estado actualizado en las consultas de categorías. \\ \hline
CU-30 & RF-INV-05 & Dueño & Usuario autenticado y autorizado para el caso de uso. & Dado un usuario autorizado, cuando registre, edite o desactive un proveedor, entonces el sistema deberá guardar el cambio y reflejar el estado actualizado en las consultas de proveedores. \\ \hline
CU-31 & RF-CFG-01 & Dueño & Usuario autenticado y autorizado para el caso de uso. & Dado un porcentaje de IVA válido y la configuración tributaria de un producto, cuando el Dueño confirme el cambio, entonces el sistema deberá aplicar la nueva configuración a operaciones posteriores y conservar el valor anterior, el nuevo valor, el actor y la fecha. \\ \hline
CU-32 & RF-BCK-01 & Superadministrador & Superadministrador autenticado y almacenamiento de respaldo disponible. & Dado un Superadministrador autenticado, cuando solicite un respaldo, entonces el sistema deberá iniciar el proceso cifrado, asignar un identificador y mostrar el estado y la fecha de finalización cuando concluya. \\ \hline
CU-33 & RF-BCK-02 & Superadministrador & Superadministrador autenticado y almacenamiento de respaldo disponible. & Dado un respaldo íntegro seleccionado, cuando el Superadministrador confirme explícitamente la restauración, entonces el sistema deberá verificar su integridad, ejecutar la restauración y registrar la operación en auditoría; si la verificación falla, no deberá restaurar. \\ \hline
CU-34 & RF-CFG-02 & Dueño & Usuario autenticado y autorizado para el caso de uso. & Dado un Dueño autenticado y una configuración válida de días y horas, cuando guarde el horario, entonces el sistema deberá actualizarlo y mostrar al Cliente únicamente el horario vigente. \\ \hline
CU-35 & RF-IA-RUT-01 & Cliente/Instructor & Actor autenticado; datos de entrada disponibles; componente de IA habilitado para evaluación. & Dado un objetivo operativo, disponibilidad semanal, historial no sensible y catálogo autorizado, cuando se solicite una recomendación, entonces el componente deberá devolver como máximo tres alternativas ordenadas de rutina y una explicación de los factores utilizados. \\ \hline
CU-36 & RF-IA-RUT-02 & Cliente & Actor autenticado; datos de entrada disponibles; componente de IA habilitado para evaluación. & Dado que el Cliente visualiza una recomendación de rutina, cuando seleccione aceptar, rechazar o solicitar alternativa, entonces el sistema deberá registrar la decisión y no deberá activar automáticamente la recomendación. \\ \hline
CU-37 & RF-IA-RUT-03 & Instructor & Actor autenticado; datos de entrada disponibles; componente de IA habilitado para evaluación. & Dado una recomendación aceptada por el Cliente, cuando el Instructor la revise, entonces el sistema deberá permitir aprobarla o rechazarla con motivo; únicamente una aprobación explícita podrá habilitar su activación. \\ \hline
CU-38 & RF-IA-ABA-01 & Dueño & Actor autenticado; datos de entrada disponibles; componente de IA habilitado para evaluación. & Dado el historial de asistencia, vigencia de membresía y frecuencia histórica de un cliente, cuando se ejecute el cálculo semanal, entonces el componente deberá asignar un nivel bajo, medio o alto y registrar la fecha del cálculo. \\ \hline
CU-39 & RF-IA-ABA-02 & Dueño/Recepcionista & Actor autenticado; datos de entrada disponibles; componente de IA habilitado para evaluación. & Dado que existen clientes clasificados con riesgo medio o alto, cuando Dueño o Recepcionista consulte las alertas, entonces el sistema deberá mostrar el nivel, la fecha y los principales factores de cada alerta sin ejecutar acciones automáticas. \\ \hline
CU-40 & RF-IA-ABA-03 & Recepcionista & Actor autenticado; datos de entrada disponibles; componente de IA habilitado para evaluación. & Dado una alerta de riesgo visible, cuando Recepción registre contacto, aplazamiento o descarte, entonces el sistema deberá guardar la acción, responsable, fecha y observación para auditoría. \\ \hline
\end{longtable}
\normalsize
\end{landscape}
